\section{Thực nghiệm và kết quả}

\subsection{Thiết lập thực nghiệm}

Hệ thống được kiểm thử trên hai cấu hình: GPU (Intel Core i7-12700H, NVIDIA RTX 3050 Ti) và CPU (Intel Core i7-1165G7, Intel Iris Xe). Camera webcam độ phân giải 640$\times$480 ở 20 FPS, YOLOv8 chạy đầu vào 640$\times$640. Dữ liệu huấn luyện: module dây an toàn sử dụng 1.200 ảnh cabin (1 lớp seat\_belt), module điện thoại sử dụng YOLOv8n tiền huấn luyện trên COCO (117.000 ảnh, 80 lớp) kết hợp fine-tuning.

\subsection{Kết quả và đánh giá}

Hệ thống có cơ chế đo thời gian suy luận (timing instrumentation) ghi log mỗi 100 frame. Bảng~\ref{tab:thresholds} liệt kê các tham số ngưỡng phát hiện của từng module.

Việc lựa chọn ngưỡng ảnh hưởng trực tiếp đến độ chính xác và tốc độ phát hiện. Ngưỡng EAR = 0,30 kết hợp với thời gian nhắm mắt tối thiểu 2 giây giúp phân biệt chớp mắt tự nhiên (100--400~ms) với buồn ngủ thực sự, giảm dương tính giả. Ngưỡng YAWN = 25~pixel và 15 khung hình liên tiếp (tương đương 0,75~giây ở 20~FPS) loại bỏ các chuyển động miệng ngắn (nói chuyện, cười). Ngưỡng tư thế đầu $|yaw| > 40^\circ$ và $pitch > 35^\circ$ được chọn dựa trên nghiên cứu về góc nhìn tự nhiên khi lái xe. Ngưỡng tin cậy 0,5 cho YOLO và FistScore $\geq 3$ (cơ chế đa số với 4/5 ngón) cân bằng giữa độ nhạy và độ đặc hiệu. Cơ chế cooldown (2--4~giây) ngăn cảnh báo trùng lặp, giảm nhiễu cho người lái.

\begin{table}[H]
\centering
\caption{Tham số ngưỡng phát hiện các module}
\fontsize{7}{8.5}\selectfont
\setlength{\tabcolsep}{2.5pt}
\begin{tabular}{llr}
\toprule
Module & Tham số & Giá trị \\
\midrule
Phát hiện nhắm mắt & EAR\_THRESHOLD & 0,30 \\
 & EAR\_MIN\_DURATION & 2 giây \\
Phát hiện ngáp & YAWN\_THRESHOLD & 25 pixel \\
 & YAWN\_CONSEC\_FRAMES & 15 khung hình \\
Tư thế đầu & Ngưỡng yaw & $>40^\circ$ \\
 & Ngưỡng pitch & $>35^\circ$ \\
Điện thoại & Confidence & $>0,5$ \\
Dây an toàn & Confidence & $>0,5$ \\
Tay lái & warning\_duration & 3 giây \\
 & FistScore & $\geq 3$ \\
\midrule
Cooldown & eye/yawn/head/phone & 2--3 giây \\
 & seatbelt/hand & 4 giây \\
\bottomrule
\end{tabular}
\label{tab:thresholds}
\end{table}


\begin{table}[H]
\centering
\caption{Mức ưu tiên và hành động xử lý cảnh báo}
\fontsize{7}{8.5}\selectfont
\setlength{\tabcolsep}{2.5pt}
\begin{tabular}{llcc}
\toprule
Module & Mức & Âm thanh & Kênh truyền \\
\midrule
Nhắm mắt & critical & nham\_mat.wav & WebSocket + MQTT \\
Điện thoại & critical & not\_phone.wav & WebSocket + MQTT \\
Dây an toàn & critical & seatbelt\_alert.wav & WebSocket + MQTT \\
Ngáp & warning & ngap\_ngu.wav & WebSocket \\
Mất tập trung & warning & dau\_quay.wav & WebSocket \\
Tay lái & warning & tay\_lai\_xe.wav & WebSocket \\
\bottomrule
\end{tabular}
\label{tab:priority}
\end{table}

Bảng~\ref{tab:inference} trình bày thời gian suy luận của từng mô-đun trên CPU với độ phân giải 640$\times$480. Pipeline xử lý tuần tự đạt tốc độ 2,16 FPS (462,37 ms/frame), trong đó các mô hình YOLOv8 chiếm phần lớn thời gian xử lý.

Độ chính xác của từng module chịu ảnh hưởng trực tiếp từ ngưỡng phát hiện. dlib đạt 95,0\% nhờ HOG+SVM phát hiện khuôn mặt chính diện ổn định, nhưng giảm khi góc nghiêng vượt $30^\circ$. YOLOv8n phát hiện điện thoại đạt precision 99,8\% ở ngưỡng conf $>0,5$ — rất cao nhờ backbone CSPDarknet, nhưng tỷ lệ phát hiện thực tế chỉ 55\% do khác biệt góc camera và điều kiện ánh sáng so với tập COCO. YOLOv8 phát hiện dây an toàn đạt 80,7\% với conf $>0,5$, thấp hơn do tập huấn luyện chỉ 1.200 ảnh. MediaPipe đạt 90,5\% nhờ mô hình TFLite nhẹ hoạt động ổn định trong cabin. Độ chính xác trung bình toàn hệ thống 89,9\% phản ánh sự cân bằng giữa độ nhạy và độ đặc hiệu ở các ngưỡng đã chọn. Về thời gian, hai mô hình YOLOv8 chiếm $\sim$71\% tổng thời gian xử lý (327 ms), cho thấy đây là nút thắt cổ chai cần tối ưu.

\begin{table}[H]
\centering
\caption{Hiệu năng suy luận của các mô-đun trong hệ thống}
\fontsize{7}{8.5}\selectfont
\setlength{\tabcolsep}{2.5pt}
\resizebox{\columnwidth}{!}{%
\begin{tabular}{lrrr}
\toprule
Mô-đun & Thời gian (ms) & FPS & Độ chính xác (\%) \\
\midrule
Phát hiện khuôn mặt và điểm mốc (dlib) & 51,96 & 19,24 & 95,0 \\
Phát hiện điện thoại (YOLOv8n-custom) & 173,23 & 5,77 & 99,8 \\
Phát hiện dây an toàn (YOLOv8-custom) & 153,92 & 6,49 & 80,7 \\
Ước lượng tư thế (MediaPipe) & 83,26 & 12,01 & 90,5 \\
Toàn bộ hệ thống & 462,37 & 2,16 & 89,9 \\
\bottomrule
\end{tabular}
}
\label{tab:inference}
\end{table}

Trên GPU, hệ thống đạt tốc độ 20--30 FPS với độ trễ 100--200 ms, trong khi trên CPU chỉ đạt 2--5 FPS với độ trễ 200--500 ms (Bảng~\ref{tab:comparison}). GPU nhanh hơn 6--10 lần, khẳng định vai trò thiết yếu của tăng tốc phần cứng trong các hệ thống DMS đa module.

\begin{table}[H]
\centering
\caption{So sánh hiệu năng CPU vs GPU}
\fontsize{7}{8.5}\selectfont
\setlength{\tabcolsep}{2.5pt}
\begin{tabular}{lrrr}
\toprule
Cấu hình & FPS & Độ trễ (ms) & Độ chính xác \\
\midrule
CPU (i7-1165G7, Iris Xe) & 2--5 & 200--500 & $\sim$85\% \\
GPU (i7-12700H, RTX 3050 Ti) & 20--30 & 100--200 & $\sim$92\% \\
\bottomrule
\end{tabular}
\label{tab:comparison}
\end{table}
